% ============================================================
%  3.2 – Demo code  (1 điểm)
%  Trạng thái: ĐÃ HOÀN THIỆN
% ============================================================

\subsection{Tổng quan demo code}

Phần này trình bày các đoạn code đại diện cho các tính năng cốt lõi đã triển khai trong
UniJob. Mã nguồn đầy đủ được lưu tại GitHub repository của nhóm.

\subsection{Xác thực người dùng -- Firebase Google OAuth}

\begin{lstlisting}[caption={auth.service.ts -- Google OAuth sign-in}]
import {
  getAuth, signInWithPopup,
  GoogleAuthProvider, onAuthStateChanged
} from 'firebase/auth';
import { doc, setDoc, getDoc, serverTimestamp } from 'firebase/firestore';
import { db } from '@/lib/firebase';

const auth = getAuth();
const provider = new GoogleAuthProvider();

export async function signInWithGoogle(): Promise<void> {
  const result = await signInWithPopup(auth, provider);
  const user = result.user;

  // Create user document if first-time login
  const userRef = doc(db, 'users', user.uid);
  const snap = await getDoc(userRef);
  if (!snap.exists()) {
    await setDoc(userRef, {
      uid: user.uid,
      displayName: user.displayName,
      email: user.email,
      photoURL: user.photoURL,
      ratingScore: null,
      skills: [],
      createdAt: serverTimestamp(),
    });
  }
}
\end{lstlisting}

\subsection{Ứng tuyển công việc -- JobDetail handleApply}

\begin{lstlisting}[caption={JobDetail.tsx -- handleApply function}]
const handleApply = async () => {
  if (!isAuthenticated || !userProfile) {
    navigate('/login');
    return;
  }
  // Guard: prevent duplicate applications
  if (hasApplied) {
    toast.error('Already applied!');
    return;
  }
  setApplying(true);
  try {
    await applyForJob(
      id!,
      userProfile.uid,
      userProfile.displayName,
      userProfile.photoURL || '',
      applyMessage
    );
    // Optimistic local state update
    setJob((prev) =>
      prev
        ? { ...prev,
            applicants: [...(prev.applicants || []), userProfile.uid] }
        : prev
    );
    toast.success('Application submitted!');
    setShowApplyForm(false);
    setApplyMessage('');
  } catch (error) {
    toast.error('Failed to apply, please try again');
  } finally {
    setApplying(false);
  }
};
\end{lstlisting}

\subsection{Thuật toán AI Job Matching -- suggest.service.ts}

\begin{lstlisting}[caption={suggest.service.ts -- Scoring algorithm}]
export function getSuggestedJobs(
  user: User,
  allJobs: Job[],
  userApplications: Application[]
): ScoredJob[] {
  const appliedJobIds = new Set(
    userApplications.map((a) => a.jobId)
  );

  return allJobs
    .filter((job) =>
      job.postedBy !== user.uid &&
      !appliedJobIds.has(job.id) &&
      job.status === 'open'
    )
    .map((job) => {
      let score = 0;
      const reasons: string[] = [];

      // Faculty match: 40 pts
      if (user.faculty && job.faculty === user.faculty) {
        score += 40;
        reasons.push('Same faculty');
      }

      // Skills/tags overlap: 30 pts
      if (user.skills?.length > 0 && job.tags?.length > 0) {
        const userSkills = new Set(
          user.skills.map((s) => s.toLowerCase())
        );
        const matched = job.tags.filter(
          (t) => userSkills.has(t.toLowerCase())
        );
        if (matched.length > 0) {
          score += Math.round(
            (matched.length / job.tags.length) * 30
          );
          reasons.push(`Skills: ${matched.join(', ')}`);
        }
      }

      // Recency: 10 pts (7-day decay)
      if (job.createdAt) {
        const age = Date.now() - job.createdAt.toDate().getTime();
        const SEVEN_DAYS = 7 * 24 * 60 * 60 * 1000;
        if (age < SEVEN_DAYS)
          score += Math.round((1 - age / SEVEN_DAYS) * 10);
      }

      // Urgent bonus: +5 pts
      if (job.isUrgent) { score += 5; }

      return { ...job, matchScore: score, matchReasons: reasons };
    })
    .sort((a, b) => b.matchScore - a.matchScore)
    .slice(0, 6);
}
\end{lstlisting}

\subsection{Xuất CV -- cv.service.ts}

\begin{lstlisting}[caption={cv.service.ts -- Canvas 2D PDF rendering (excerpt)}]
export async function generateCVCertificate(
  user: User,
  completedJobs: Job[],
  ratings: Rating[],
  options: CVExportOptions = { showDetailedRatings: false }
): Promise<void> {
  const canvas = document.createElement('canvas');
  const SCALE = 3; // 3x resolution for sharp PDF
  canvas.width  = 794 * SCALE;  // A4 width at 96dpi
  canvas.height = 1123 * SCALE; // A4 height at 96dpi
  const ctx = canvas.getContext('2d')!;
  ctx.scale(SCALE, SCALE);

  // White background
  ctx.fillStyle = '#ffffff';
  ctx.fillRect(0, 0, 794, 1123);

  // Green header bar
  ctx.fillStyle = '#10b981'; // emerald-500
  ctx.fillRect(0, 0, 794, 14);

  // Certificate title
  ctx.font = 'bold 22px "Segoe UI", sans-serif';
  ctx.fillStyle = '#1f2937';
  ctx.textAlign = 'center';
  ctx.fillText('EXPERIENCE CERTIFICATE', 794 / 2, 120);

  // ... (stat boxes, job list, footer)

  // Embed canvas into jsPDF
  const imgData = canvas.toDataURL('image/png');
  const pdf = new jsPDF('p', 'mm', [210, 297]);
  pdf.addImage(imgData, 'PNG', 0, 0, 210, 297);
  pdf.save(`certificate_${user.displayName}.pdf`);
}
\end{lstlisting}

\textit{Lý do dùng Canvas 2D thay vì html2canvas:} Thư viện html2canvas gặp lỗi parse
màu \texttt{oklch()} của Tailwind CSS v4. Canvas 2D API vẽ trực tiếp bằng hex color,
không phụ thuộc vào CSS parser của browser -- đảm bảo export PDF ổn định trên mọi
trình duyệt.

\subsection{Tìm kiếm -- Client-side search filter}

\begin{lstlisting}[caption={jobStore.ts -- Client-side search query filter}]
// After fetching from Firestore, filter by searchQuery
const { searchQuery } = get().filters;
if (searchQuery?.trim()) {
  const q = searchQuery.toLowerCase();
  fetchedJobs = fetchedJobs.filter(
    (job) =>
      job.title.toLowerCase().includes(q) ||
      job.description?.toLowerCase().includes(q)
  );
}
\end{lstlisting}

\textit{Ghi chú:} Firestore không hỗ trợ full-text search nội địa. Giải pháp lọc
client-side đủ dùng cho quy mô demo ($<$500 records). Trong production, có thể tích hợp
Algolia hoặc Typesense.

\subsection{Hệ thống đánh giá uy tín -- rating.service.ts}

\begin{lstlisting}[caption={rating.service.ts -- submitRating và cập nhật điểm uy tín tự động}]
export async function submitRating(
  jobId: string,
  fromUserId: string,
  toUserId: string,
  score: number,        // 1-5
  comment: string,
  type: 'poster-to-worker' | 'worker-to-poster'
): Promise<string> {
  // Luu rating vao Firestore
  const docRef = await addDoc(ratingsCollection, {
    jobId, fromUserId, toUserId,
    score, comment, type,
    createdAt: serverTimestamp(),
  });

  // Tinh lai diem uy tin trung binh cho nguoi nhan rating
  const userRatings = await getRatingsByUser(toUserId);
  const avg = userRatings.reduce((s, r) => s + r.score, 0)
              / userRatings.length;

  await updateDoc(doc(db, 'users', toUserId), {
    ratingScore:  Math.round(avg * 10) / 10,
    totalRatings: userRatings.length,
  });
  return docRef.id;
}
\end{lstlisting}

\subsection{Xác nhận hoàn thành công việc 2 chiều -- confirmCompletion}

\begin{lstlisting}[caption={rating.service.ts -- Xác nhận hoàn thành 2 chiều (poster + worker)}]
export async function confirmCompletion(
  completionId: string,
  role: 'worker' | 'poster'
): Promise<void> {
  const ref  = doc(db, 'jobCompletions', completionId);
  const snap = await getDoc(ref);
  if (!snap.exists()) return;

  const data  = snap.data() as JobCompletion;
  const field = role === 'worker' ? 'workerConfirmed' : 'posterConfirmed';
  const other = role === 'worker' ? data.posterConfirmed : data.workerConfirmed;

  const updates: any = { [field]: true };
  if (other) {
    // Ca hai ben da xac nhan -- hoan thanh chinh thuc
    updates.status      = 'confirmed';
    updates.completedAt = serverTimestamp();
  }
  await updateDoc(ref, updates);

  if (other) {
    // Ghi nhan workHistory va dong task
    const job = await getJobById(data.jobId);
    if (job?.assignedTo?.[0]) {
      await recordWorkHistory(job.postedBy, job.assignedTo[0], job.id);
      await updateJob(job.id, { status: 'completed' });
    }
  }
}
\end{lstlisting}

\textit{Cơ chế 2-step confirmation:} Cả poster lẫn worker phải nhấn xác nhận
(\texttt{posterConfirmed = true} \textbf{và} \texttt{workerConfirmed = true}) thì hệ
thống mới ghi nhận \texttt{workHistory}, mở khóa chức năng đánh giá, và cập nhật
trạng thái task thành \texttt{completed}. Cơ chế này tránh gian lận đánh giá
và đảm bảo tính công bằng cho cả hai bên.

\subsection{State Management -- Zustand authStore}

\begin{lstlisting}[caption={authStore.ts -- Global auth state with Zustand}]
interface AuthState {
  user: FirebaseUser | null;
  userProfile: User | null;
  isAuthenticated: boolean;
  isLoading: boolean;
  initialize: () => Unsubscribe;
}

export const useAuthStore = create<AuthState>((set) => ({
  user: null,
  userProfile: null,
  isAuthenticated: false,
  isLoading: true,

  initialize: () => {
    return onAuthStateChanged(auth, async (firebaseUser) => {
      if (firebaseUser) {
        const profile = await getUserById(firebaseUser.uid);
        set({
          user: firebaseUser,
          userProfile: profile,
          isAuthenticated: true,
          isLoading: false,
        });
      } else {
        set({
          user: null,
          userProfile: null,
          isAuthenticated: false,
          isLoading: false,
        });
      }
    });
  },
}));
\end{lstlisting}
